\begin{proofbox}
\textbf{方法一：一次平方，直接证明方程等价（优先）。}
按定义建系，$F(p/2,0)$，$\ell:x=-p/2$，$p>0$。
设 $P(x,y)$，则 $H(-p/2,y)$。两边都是非负距离，所以平方是等价变形：
\[
\begin{aligned}
PF=PH
&\iff \sqrt{(x-p/2)^2+y^2}=|x+p/2|\\
&\iff (x-p/2)^2+y^2=(x+p/2)^2\\
&\iff \boxed{y^2=2px}.
\end{aligned}
\]
标准方程还给出 $x\ge0$，故 $PH=x+p/2\ge p/2>0$，
既说明没有点在准线上，也给出焦半径的直接计算式 $PF=x+p/2$。

\textbf{方法二：垂直平分线构造（只用初中几何）。}
在准线上任取 $H$，过 $H$ 作 $\ell$ 的垂线，再作 $FH$ 的垂直平分线；
两线交于 $P$。因为 $P$ 在垂直平分线上，所以 $PF=PH$，正好满足定义。
反过来，每个满足定义的点 $P$，取它在准线上的垂足 $H$，也必在 $FH$ 的垂直平分线上。
因此构造没有漏点或添点。
\end{proofbox}
\clearpage
\begin{proofbox}
\begin{center}\begin{tikzpicture}[font=\small,>=stealth]
\draw[black!30,rounded corners=2pt] (0,0) rectangle (14.4,8.2);
\node[anchor=west,font=\small\bfseries] at (.5,7.72) {几何构造：先在准线上取 $H$，再作 $FH$ 的垂直平分线};
\begin{scope}[shift={(5.1,3.4)},scale=.85]
\draw[violet,thick] (-1,-2.5)--(-1,3.5) node[above] {$\ell$};
\draw[blue!30,domain=-2.5:3.5,samples=121,variable=\t] plot ({\t*\t/4},\t);
\draw[black!50] (-1,2)--(4,2);
\draw[orange!85!black,very thick] (1,0)--(-1,2);
\draw[blue,very thick] (-1.5,-.5)--(2.5,3.5);
\draw[orange!85!black,dashed] (1,0)--(1,2)--(-1,2);
\draw[black!60] (-1,1.75)--(-.75,1.75)--(-.75,2);
\fill (1,0) circle (2pt) node[below] {$F$};
\fill (-1,2) circle (2pt) node[left] {$H$};
\fill (0,1) circle (2pt) node[left=7pt] {$M$};
\fill (1,2) circle (2pt) node[right=8pt] {$P$};

\end{scope}
\node[anchor=west] at (8.4,6.6) {$FH$ 的垂直平分线};
\draw[black!50,->] (8.2,6.5)--(7.225,6.375);

\node[anchor=west] at (.7,.55) {$PH\perp\ell$，$MF=MH$；交点 $P$ 满足 $PF=PH$，移动 $H$ 得到抛物线};
\end{tikzpicture}\end{center}

图中 $H=(-1,2),F=(1,0)$，中点 $M=(0,1)$，垂直平分线为 $y=x+1$。
它与 $y=2$ 交于 $P=(1,2)$。
一般地，$FH$ 的垂直平分线不会平行于过 $H$ 的准线垂线：否则 $FH$ 就平行于准线，与 $F\notin\ell$ 矛盾。
所以每个 $H$ 都唯一确定一个轨迹点。
\end{proofbox}

% ===== C038 补充：新增思路，不改变原两种证明及 TikZ =====
\clearpage
\begin{proofbox}
\textbf{方法三：先代入最简单的数字，再认出平方差（忘记公式时最快恢复）。}
只看 $p=2$，此时 $F(1,0)$，$\ell:x=-1$，设 $P(x,y)$。
\[
PF^2=(x-1)^2+y^2,\qquad PH^2=(x+1)^2.
\]
因为 $PF=PH$，所以直接令\textbf{平方相等}，而不用计算两个平方根：
\[
\begin{aligned}
(x-1)^2+y^2&=(x+1)^2,\\
y^2&=(x+1)^2-(x-1)^2\\
&=\big[(x+1)-(x-1)\big]\big[(x+1)+(x-1)\big]\\
&=2\cdot2x=\boxed{4x}.
\end{aligned}
\]
\textbf{整个诀窍：两个平方一相减，$x^2$ 自动消掉。}

换回一般情况，只须把上面的 $1$ 换成 $p/2$：
\[
\boxed{\ y^2=(x+p/2)^2-(x-p/2)^2
=\underbrace{p}_{\text{两式之差}}\cdot\underbrace{2x}_{\text{两式之和}}=2px\ }.
\]
这是\textbf{考场速写法}：写出左右两个平方，直接用平方差，少一次展开。
\end{proofbox}
\clearpage
\begin{proofbox}
\textbf{方法四：只用初中几何，先确定顶点和对称轴。}
令 $D$ 为焦点 $F$ 在准线上的垂足。
\begin{enumerate}[leftmargin=*]
\item \textbf{对称轴为什么是 $FD$？} 关于直线 $FD$ 对称时，焦点 $F$ 不动，准线 $\ell$ 也与自身重合。
因此，点 $P$ 满足的等距条件在对称后仍成立；轨迹关于 $FD$ 对称。
\item \textbf{顶点为什么在中点？} 在 $FD$ 上找一个满足定义的点 $O$。
因 $O$ 到准线的垂距为 $OD$，于是 $OF=OD$，所以 $O$ 是 $FD$ 的中点。
\item \textbf{为什么开口背离准线？} 按定义求出的标准方程是 $y^2=2px$。
因为 $p>0$ 且 $y^2\ge0$，所以 $x\ge0$；轨迹位于过顶点、平行于准线的直线的\textbf{焦点一侧}。
\end{enumerate}
前两步完全不需要坐标，第三步只需一个平方数非负的事实。
\end{proofbox}
\clearpage
\begin{proofbox}
\textbf{方法五：反过来验证方程，一步确认没有多余点。}
假设 $P(x,y)$ 满足 $y^2=2px$，$p>0$。
由于 $x\ge0$，有 $x+p/2>0$。于是
\[
\begin{aligned}
PF^2&=(x-p/2)^2+y^2\\
&=(x-p/2)^2+2px\\
&=(x+p/2)^2=PH^2.
\end{aligned}
\]
两边是非负的长度，故 $PF=PH$，点 $P$ 必满足定义。

\textbf{意义：}原有“由等距推出方程”说明\textbf{没有漏掉满足定义的点}；
这里“由方程推出等距”说明\textbf{方程没有引入多余的点}。
两个方向合起来，才是严格的“当且仅当”。

\textbf{顺手得到的结论：}
\[
\boxed{PF=PH=x+\frac p2\ge\frac p2}.
\]
所以从抛物线上任意一点到焦点的距离，最小值是 $p/2$，
并且只有在顶点 $O=(0,0)$ 处取到。
\end{proofbox}
